% ============================================================
% CHƯƠNG 2: CƠ SỞ LÝ THUYẾT VÀ NGHIÊN CỨU LIÊN QUAN
% MỤC ĐÍCH: Đáp ứng Câu 3 (nghiên cứu tổng quan ≥3 vấn đề, ≥3 nghiên cứu)
% PHÂN BỔ: ~10–12 trang
% ============================================================
\chapter{Cơ sở lý thuyết và nghiên cứu liên quan}
\label{chap:co-so-ly-thuyet}

Chương này tập hợp nền tảng lý thuyết cho ELDAS theo bốn mảng: (i)~các thuật
toán lập lịch trong trung tâm dữ liệu; (ii)~mô hình tiêu thụ năng lượng;
(iii)~học tăng cường đa mục tiêu (MORL); và (iv)~các công trình nghiên cứu
liên quan. Phần trình bày được tiết chế ở mức vừa đủ để người đọc nắm được
các quyết định thiết kế ở Chương~\ref{chap:gioi-thieu} cũng như phần hiện
thực ở các chương sau, đồng thời định vị được đóng góp của đề tài trong
tương quan với các nghiên cứu đã có.

% ============================================================
\section{Lập lịch trong trung tâm dữ liệu}
\label{sec:lap-lich-dc}

\subsection{Các thuật toán heuristic truyền thống}
Bài toán lập lịch tài nguyên đã được nghiên cứu từ rất sớm, dưới dạng
\textit{Bin Packing}: xếp $n$ vật phẩm có kích thước cho trước vào số thùng
nhỏ nhất, mỗi thùng có dung lượng cố định~\cite{garey1979computers}. Một số
\textit{heuristic} kinh điển vẫn được sử dụng phổ biến cho đến nay:

\begin{itemize}
  \item \textbf{First-Fit (FF).} Duyệt danh sách host theo thứ tự cố định,
        đặt tác vụ lên host đầu tiên còn đủ tài nguyên. Cách làm đơn giản,
        độ phức tạp $O(n \cdot m)$, nhưng dễ gây phân mảnh.
  \item \textbf{Best-Fit (BF).} Đặt tác vụ lên host \textit{vừa khít nhất}
        --- tức host có tài nguyên dư còn lại nhỏ nhất sau khi đặt. Giảm
        phân mảnh nhưng tốn chi phí sắp xếp.
  \item \textbf{Worst-Fit (WF).} Ngược lại với BF: chọn host còn dư nhiều
        nhất để cân bằng tải, đổi lại làm tăng số host phải bật.
  \item \textbf{Round-Robin (RR).} Phân phối tuần tự, bảo đảm công bằng
        nhưng bỏ qua sự khác biệt về dung lượng giữa các host.
\end{itemize}

Điểm chung của các heuristic này là \textbf{đơn mục tiêu}
(single-objective): hoặc tối thiểu số host, hoặc cân bằng tải, không tính
đến năng lượng hay SLA. Đây là hạn chế đáng kể khi áp dụng cho các
datacenter AI hiện đại.

\subsection{Meta-heuristic: Ant Colony Optimization (ACO)}
\label{ssec:aco}

Để vượt qua hai giới hạn cố hữu của các heuristic ``nhìn một bước''
(FF, BF) là \emph{đơn mục tiêu} và \emph{dễ kẹt ở cực tiểu địa
phương}, nhiều nghiên cứu chuyển sang nhóm \textbf{meta-heuristic};
trong đó \textit{Ant Colony Optimization} (ACO) là đại diện được khảo
sát rộng rãi nhất cho bài toán lập lịch tài
nguyên~\cite{dorigo1996ant}.

\paragraph{Ý tưởng cốt lõi.} ACO mô phỏng cách đàn kiến tìm đường ngắn
nhất tới nguồn thức ăn dựa trên \emph{pheromone trail}: mỗi con kiến đi
qua một đoạn đường sẽ để lại một lượng pheromone; đoạn nào tích luỹ
nhiều pheromone hơn sẽ được các con kiến sau ưu tiên chọn, tạo thành
vòng phản hồi dương và dần hội tụ về đường tối ưu. Khi ánh xạ sang bài
toán scheduling, mỗi ``con kiến'' xây một schedule hoàn chỉnh bằng
cách chọn host cho từng task theo xác suất tỉ lệ thuận với lượng
pheromone đã tích luỹ và chất lượng ước tính của cặp
\textit{(task, host)}. Sau mỗi vòng, pheromone bay hơi một phần rồi
được củng cố theo mức độ tốt của schedule vừa tìm ra, hình thành ``bộ
nhớ tập thể'' để định hướng cho các kiến ở vòng sau. Liu và cộng sự áp
dụng biến thể \textit{Energy-aware ACO} cho bài toán VM placement và
giảm được $9$--$12\%$ năng lượng so với First-Fit Decreasing trên cùng
một bộ workload trace~\cite{liu2014adaptive}.

\paragraph{Ưu điểm.} (i)~khả năng \emph{thoát khỏi cực tiểu địa phương}
nhờ tính ngẫu nhiên trong lựa chọn của từng kiến; (ii)~hỗ trợ
multi-objective một cách tự nhiên thông qua hàm chi phí tổng hợp;
(iii)~không cần gradient, có thể hoạt động trên không gian rời rạc
rộng.

\paragraph{Hạn chế đối với ELDAS.} Trong phạm vi đồ án này, ACO bộc lộ
ba nhược điểm đáng kể:
\begin{enumerate}
  \item \textbf{Chi phí đánh giá rất lớn.} ACO thuộc nhóm
        population-based: mỗi iteration cần $K$ kiến
        ($K \approx 20$--$50$) $\times$ một episode mô phỏng đầy đủ.
        Với throughput đo được khoảng $65$\,step/s
        (§\ref{ssec:res-throughput}) và $\sim 2\,000$ task/episode,
        một iteration ACO tiêu tốn khoảng $30$\,s $\times 30$ kiến
        $\approx 15$ phút; $100$ iteration hội tụ tương đương
        $\sim 25$ giờ \emph{cho một bộ siêu tham số duy nhất}.
  \item \textbf{Khó khái quát hoá.} Ma trận pheromone $\tau_{ij}$ được
        học riêng cho một workload cụ thể; khi gặp workload mới, phải
        chạy lại từ đầu. Ngược lại, một chính sách RL đã huấn luyện
        có khả năng khái quát hoá nhờ hàm xấp xỉ.
  \item \textbf{Nhiều siêu tham số.} ACO có khá nhiều tham số cần
        chỉnh ($\alpha, \beta, \rho, Q, K$, chiến lược elitism, \dots),
        khiến khả năng tái lập kết quả thấp và việc so sánh công bằng
        giữa các nghiên cứu trở nên khó khăn.
\end{enumerate}

\textbf{So sánh với PPO (định hướng của ELDAS).} PPO khắc phục được cả
ba điểm yếu kể trên: (i)~tuy là on-policy, PPO cập nhật ngay sau
\emph{mỗi episode} thay vì sau \emph{cả một quần thể episode}, tiết
kiệm khoảng $K$ lần thời gian mô phỏng; (ii)~khái quát hoá qua mạng
nơ-ron --- một chính sách có thể dùng cho nhiều workload; (iii)~chỉ có
một vài siêu tham số chính (learning rate, clip $\epsilon$, GAE
$\lambda$). Do đó, ACO chỉ được nhắc đến như \emph{tham chiếu so sánh}
trong phần khảo sát, không được hiện thực trong ELDAS.

\subsection{Kubernetes Scheduler}
Kubernetes (K8s) là hệ điều phối container phổ biến nhất hiện
nay~\cite{kubernetes2024scheduler}. Bộ lập lịch mặc định
(\textit{kube-scheduler}) hoạt động theo \textit{pipeline} ba bước cho mỗi
\textit{pod}:

\begin{enumerate}
  \item \textbf{Filter (Predicate).} Loại bỏ các node không đủ điều kiện. Một
        số bộ lọc tiêu biểu: \texttt{PodFitsResources} (đủ CPU/RAM/GPU),
        \texttt{NodeAffinity}, \texttt{Taints/Tolerations}.
  \item \textbf{Score (Priority).} Chấm điểm từng node còn lại trên thang
        $[0,10]$ bằng các hàm tính điểm. Hai hàm được dùng trong ELDAS để mô
        phỏng K8s:
        \[
          \mathrm{LRP}(h_j) = 10 \cdot
          \frac{(c^{\mathrm{cpu}}_j - u^{\mathrm{cpu}}_j)
              + (c^{\mathrm{ram}}_j - u^{\mathrm{ram}}_j)}
              {2\,(c^{\mathrm{cpu}}_j + c^{\mathrm{ram}}_j)},
        \]
        gọi là \textit{LeastRequestedPriority} -- ưu tiên host có nhiều tài
        nguyên trống (cân bằng tải). Hàm thứ hai
        \textit{BalancedResourceAllocation} ưu tiên host có tỉ lệ CPU/RAM sử
        dụng đồng đều, nhằm tránh nghẽn cổ chai một chiều tài nguyên.
  \item \textbf{Bind.} Gắn pod vào node có điểm số cao nhất.
\end{enumerate}

\textbf{Hạn chế.} Bộ lập lịch của K8s không có hàm chấm điểm theo năng
lượng, và bản thân nó không hỗ trợ tối ưu đa mục tiêu một cách nguyên thủy.
Đây cũng chính là khoảng trống mà các phương pháp dựa trên học tăng cường
có thể khai thác.

\subsection{Lập lịch dựa trên học tăng cường: Deep Q-Network (DQN)}
\label{ssec:dqn}

Kết quả đột phá của Mnih và cộng sự trên bộ trò chơi Atari năm
$2015$~\cite{mnih2015human} đã đưa \textit{Deep Q-Network} (DQN) trở
thành một trong những thuật toán RL được khảo sát phổ biến nhất cho
bài toán scheduling trong giai đoạn sau đó. DQN học hàm giá trị
\textbf{Q-function} $Q(s, a)$ --- ước lượng tổng reward kỳ vọng khi
thực hiện hành động $a$ từ trạng thái $s$. Thay vì tính trực tiếp,
DQN dùng một mạng nơ-ron để xấp xỉ $Q$ và cập nhật theo nguyên lý
Bellman: phần thưởng hiện tại cộng với giá trị tốt nhất có thể đạt
được ở bước kế tiếp. Hai kỹ thuật then chốt giúp quá trình học ổn
định là \textbf{target network} (một bản sao mạng được đóng băng định
kỳ để tránh hiện tượng ``mục tiêu di động'') và \textbf{replay buffer}
(lưu lại kinh nghiệm rồi lấy mẫu ngẫu nhiên, phá vỡ tương quan thời
gian giữa các bước liên tiếp).

\paragraph{Ưu điểm.} (i)~\textbf{Tiết kiệm mẫu hơn so với nhóm
on-policy}: cùng một mẫu kinh nghiệm có thể được tái sử dụng nhiều
lần qua replay buffer --- một lợi thế rõ rệt khi mỗi episode mô phỏng
CloudSim đều tốn kém; (ii)~\textbf{phù hợp tự nhiên với không gian
hành động rời rạc}: phép $\arg\max_a Q$ là toán tử kiểm soát mặc
định, ăn khớp với $H = 10$ lựa chọn host; (iii)~biến thể
\textit{Rainbow DQN} (kết hợp PER, Double DQN, Dueling, Distributional)
hiện đang giữ SOTA trên nhiều benchmark.

\paragraph{Hạn chế đối với ELDAS.} Có bốn lý do khiến DQN \emph{không}
phải lựa chọn tối ưu cho phạm vi Phase 2:
\begin{enumerate}
  \item \textbf{Reward lệch biên độ dẫn tới ước lượng $Q$ thiên dương.}
        Reward energy trong ELDAS luôn âm và có biên độ vào khoảng
        $10^3$\,Ws, trong khi reward SLA thưa hơn và biên độ phụ thuộc
        vào mức vi phạm. DQN chọn hành động dựa trên $\arg\max_a Q$,
        nên dễ bị \emph{thiên dương} (overestimation) khi phân phối
        $Q$ có hai thành phần chênh nhau tới hai bậc độ lớn. PPO xử lý
        được vấn đề này bằng \textit{advantage normalisation} ---
        chuẩn hoá $A_t$ về trung bình 0, độ lệch chuẩn 1 trước mỗi
        bước cập nhật.
  \item \textbf{Action masking phức tạp hơn.} Các host không đủ tài
        nguyên cần bị cấm tuyệt đối. Trong DQN, điều này phải được
        thực thi ở hai vị trí --- bước chọn hành động theo
        $\epsilon$-greedy và bước $\max_{a'}$ trong Bellman backup ---
        với nhiều biến thể chưa được chuẩn hoá trong literature.
        Trong khi đó, với PPO chỉ cần mask một lần tại policy head, và
        thư viện \texttt{sb3-contrib} đã có sẵn \texttt{MaskablePPO}.
  \item \textbf{Không hỗ trợ vector reward một cách nguyên thuỷ.}
        Q-function trả về vô hướng cho mỗi cặp $(s, a)$; muốn dùng
        cho MORL thì phải \emph{scalarize trước khi học}, đồng nghĩa
        với việc phải huấn luyện lại từ đầu cho từng trọng số
        $\mathbf{w}$ --- mức chi phí ngang bằng với linear
        scalarization của PPO. Các biến thể MORL chuyên dụng như
        \textit{Envelope-Q} hay \textit{Pareto-DQN} có tồn tại, nhưng
        phức tạp hơn, chưa có triển khai sẵn trên SB3 và do đó nằm
        ngoài phạm vi Phase 2.
  \item \textbf{Mô hình off-policy + replay buffer khó song song hoá.}
        Phase 2 dự kiến chạy weight sweep ($K = 9$) qua
        \texttt{SubprocVecEnv} để tận dụng nhiều core CPU; PPO
        on-policy ăn khớp tự nhiên với kiến trúc này, trong khi DQN
        đòi hỏi các kiến trúc phức tạp hơn nhiều như Ape-X hay R2D2.
\end{enumerate}

\textbf{So sánh với PPO (định hướng của ELDAS).} Trước phổ reward lệch
biên độ, PPO ổn định hơn DQN nhờ hai cơ chế mà DQN không có tương
đương: \emph{advantage normalisation} --- triệt tiêu chênh lệch biên
độ trước mỗi lần cập nhật --- và \emph{trust region thông qua
clipping} --- giới hạn mức thay đổi của chính sách ở mỗi bước, tránh
hiện tượng diverge khi vi phạm SLA bất ngờ tăng vọt. Phần chi tiết về
PPO được trình bày ở §\ref{sec:morl}.

\begin{table}[H]
  \centering
  \caption{So sánh các thuật toán lập lịch truyền thống.}
  \label{tab:scheduling-comparison}
  \small
  \begin{tabular}{lccc}
    \toprule
    \textbf{Thuật toán} & \textbf{Mục tiêu} & \textbf{Đa mục tiêu?}
                       & \textbf{Năng lượng?} \\
    \midrule
    First-Fit            & Min \#host        & Không & Không \\
    Best-Fit             & Min phân mảnh     & Không & Không \\
    Round-Robin          & Cân bằng tải      & Không & Không \\
    K8s LRP              & Cân bằng tải      & Không & Không \\
    NSGA-II              & Energy + QoS      & Có    & Có    \\
    DRL (DeepRM)         & Slowdown          & Không & Không \\
    \textbf{MORL (ELDAS)}& \textbf{Energy + SLA} & \textbf{Có} & \textbf{Có} \\
    \bottomrule
  \end{tabular}
\end{table}

% ============================================================
\section{Mô hình tiêu thụ năng lượng datacenter}
\label{sec:nang-luong}

\subsection{Mô hình tuyến tính SPECpower}
Phép đo thực nghiệm trên hơn $200$ máy chủ thương mại do tổ chức SPEC công bố
(SPECpower\_ssj2008) cho thấy công suất tiêu thụ của một máy chủ phụ thuộc
gần như tuyến tính vào mức sử dụng CPU~\cite{beloglazov2012energy}. Mô hình
đơn giản và được dùng rộng rãi nhất là:
\begin{equation}
  P(U) \;=\; P_{\mathrm{idle}} \;+\; U \cdot \big(P_{\mathrm{max}} -
                                                  P_{\mathrm{idle}}\big),
  \label{eq:linear-power}
\end{equation}
trong đó:
\begin{itemize}
  \item $U \in [0,1]$ -- mức sử dụng CPU trung bình của host;
  \item $P_{\mathrm{idle}}$ -- công suất ở trạng thái rảnh ($U=0$);
  \item $P_{\mathrm{max}}$ -- công suất ở trạng thái tải tối đa ($U=1$).
\end{itemize}

Có một đặc điểm đáng chú ý: $P_{\mathrm{idle}}$ thường chiếm tới
$50$--$70\%$ của $P_{\mathrm{max}}$. Nói cách khác, một host ở trạng thái
\textit{rảnh} vẫn tiêu tốn hơn phân nửa công suất so với khi chạy hết tải.
Đây cũng là cơ sở để các thuật toán \textit{consolidation} ưu tiên gom tác
vụ về ít host nhất có thể rồi tắt số host còn lại.

\subsection{Tính tổng năng lượng theo thời gian rời rạc}
Mô phỏng trong ELDAS chạy theo \textbf{thời gian rời rạc} với bước
$\Delta t = 1\,\mathrm{s}$. Tổng năng lượng tiêu thụ của một host trong
$K$ bước được tính bằng phép cộng dồn:
\begin{equation}
  E_j \;=\; \sum_{k=1}^{K} P_j\!\left(U_j^{(k)}\right) \cdot \Delta t,
  \label{eq:energy-discrete}
\end{equation}
và năng lượng toàn cụm là $E = \sum_{j=1}^{m} E_j$. Cách rời rạc hoá này
ánh xạ trực tiếp xuống \textit{event loop} của CloudSim Plus
(xem Chương~4), nên không cần đến phép tích phân giải tích.

\subsection{Mở rộng cho GPU}
CloudSim Plus phiên bản 8.5.7 chưa có khái niệm GPU một cách nguyên thủy.
ELDAS bổ sung thêm một lớp theo dõi GPU riêng, dùng chung dạng công thức
tuyến tính với hai mức công suất $P^{\mathrm{gpu}}_{\mathrm{idle}}$ và
$P^{\mathrm{gpu}}_{\mathrm{max}}$ lấy từ tài liệu kỹ thuật của NVIDIA
A100~\cite{nvidia2023a100}.

\subsection{Power Usage Effectiveness (PUE)}
Để phản ánh chi phí năng lượng \textit{vận hành} (làm mát, chiếu sáng, mạng),
ngành công nghiệp dùng chỉ số PUE:
\begin{equation}
  \mathrm{PUE} \;=\; \frac{E_{\mathrm{total}}}{E_{\mathrm{IT}}},
  \label{eq:pue}
\end{equation}
với $E_{\mathrm{IT}}$ là năng lượng tiêu thụ bởi thiết bị tính toán và
$E_{\mathrm{total}}$ là tổng năng lượng đầu vào của datacenter. PUE lý tưởng
là $1{,}0$; trung bình toàn ngành 2023 vào khoảng $1{,}58$~\cite{uptime2023pue}.
ELDAS sử dụng giá trị mặc định $\mathrm{PUE} = 1{,}5$ để quy đổi $E_{\mathrm{IT}}$
sang chi phí thực tế khi báo cáo.

\subsection{Lý do chọn mô hình tuyến tính cho Phase 1}
Việc dừng lại ở mô hình tuyến tính trong Phase 1 dựa trên ba lý do:
\begin{enumerate}
  \item \textbf{Đủ chính xác cho mục đích so sánh tương đối.} Sai số
        trung bình so với phép đo thực nghiệm chỉ vào khoảng
        $5$--$7\%$~\cite{beloglazov2012energy} --- đủ để xếp hạng các
        scheduler.
  \item \textbf{Đơn giản và dễ tái lập.} Mỗi loại thiết bị chỉ cần hai
        tham số, vừa dễ tinh chỉnh, vừa minh bạch khi viết lại trong
        báo cáo.
  \item \textbf{Tránh phụ thuộc vào DVFS.} Các mô hình phi tuyến (bậc
        ba, logarit) và cơ chế \textit{Dynamic Voltage and Frequency
        Scaling} (DVFS) cần dữ liệu tần số/điện áp theo thời gian
        thực --- Alibaba trace không cung cấp các trường này, do đó
        phần mở rộng được giữ lại cho Phase 2 hoặc cho hướng nghiên
        cứu sau.
\end{enumerate}

% ============================================================
\section{Học tăng cường đa mục tiêu (MORL)}
\label{sec:morl}

\subsection{Markov Decision Process (MDP)}
Học tăng cường (\textit{Reinforcement Learning} --- RL) mô hình hoá bài
toán ra quyết định tuần tự dưới dạng \textit{Markov Decision
Process}~\cite{sutton2018reinforcement}. Ở mỗi bước, \textbf{agent} quan
sát \textbf{trạng thái} của môi trường, chọn một \textbf{hành động}, nhận
về một \textbf{phần thưởng} rồi chuyển sang trạng thái mới; quá trình
này được lặp lại cho đến khi episode kết thúc. Mục tiêu của agent là học
một \textbf{chính sách} --- một ánh xạ từ trạng thái sang hành động ---
sao cho tổng phần thưởng tích luỹ (có chiết khấu theo thời gian) đạt giá
trị lớn nhất.

\subsection{Multi-Objective MDP (MOMDP)}
Khi xuất hiện \textbf{nhiều mục tiêu xung đột}, phần thưởng không còn là
một số vô hướng đơn lẻ mà trở thành một \textbf{vector phần
thưởng}~\cite{roijers2013survey, hayes2022practical}. Trong ELDAS, ở
mỗi bước agent nhận đồng thời hai phần thưởng riêng biệt:
$R^{\mathrm{energy}}$ --- mức năng lượng tiêu thụ trong bước đó, luôn
mang giá trị âm --- và $R^{\mathrm{SLA}}$ --- mức vi phạm deadline,
bằng 0 nếu tác vụ xong đúng hạn.

\subsection{Pareto Optimality}
Do không thể cộng trực tiếp hai mục tiêu khác đơn vị (Joule và giây
trễ), khái niệm tối ưu phải được mở rộng thành \textbf{Pareto
optimality} (đã được định nghĩa ở §1.2). Trong không gian giá trị kỳ
vọng của vector phần thưởng, tập các chính sách không bị thống trị tạo
thành \textit{Pareto front} $\mathcal{P}^*$.

\textbf{Vì sao cần vector reward thay vì giá trị vô hướng?} Nếu cộng gộp
ngay từ đầu $R = w_1 R_1 + w_2 R_2$ với một $\mathbf{w}$ cố định, agent
chỉ học được \emph{một điểm} duy nhất trên Pareto front. Giữ phần
thưởng ở dạng vector mang lại hai lợi ích: (i)~có thể quét trọng số
sau khi huấn luyện để dựng lại nhiều điểm trên Pareto front; (ii)~có
thể áp dụng các thuật toán MORL chuyên biệt như \textit{Envelope
Q-Learning} hay \textit{Pareto-DQN} trong các nghiên cứu kế tiếp.


\subsection{Proximal Policy Optimization (PPO)}
PPO~\cite{schulman2017proximal} hiện là một trong những thuật toán
\textit{policy gradient} ổn định nhất, và cũng là lựa chọn cho Phase 2
của ELDAS. Ý tưởng chính của PPO khá đơn giản: ở mỗi vòng cập nhật, giữ
độ \emph{lệch} của chính sách mới so với chính sách cũ trong một
\textit{vùng tin cậy}. Thay vì đi giải một bài toán ràng buộc phức tạp,
PPO áp dụng cơ chế \textbf{cắt tỉ số} (clipping): tỉ số xác suất giữa
hai chính sách bị giới hạn trong khoảng $[1-\varepsilon,\, 1+\varepsilon]$
(thường lấy $\varepsilon = 0{,}2$), nhờ đó mỗi bước cập nhật không thể
kéo chính sách thay đổi quá lớn và gây diverge.

\textbf{Vì sao chọn PPO?} (i)~ổn định nhất trong nhóm policy gradient;
(ii)~ít siêu tham số, dễ tinh chỉnh; (iii)~được hỗ trợ tốt trong
\texttt{stable-baselines3}; (iv)~đã chứng minh hiệu quả trên các bài
toán scheduling tương
tự~\cite{mao2016resource, decima2019}.

\subsection{Linear Scalarization và quét trọng số}
Cách đơn giản nhất để tối ưu một vector phần thưởng là \textbf{cộng có
trọng số}:
\begin{equation}
  R^{\mathrm{scalar}}_t(\mathbf{w}) \;=\;
    w \cdot R^{\mathrm{energy}}_t \;+\; (1 - w) \cdot R^{\mathrm{SLA}}_t,
  \qquad w \in [0,1].
  \label{eq:linear-scalar}
\end{equation}
Mỗi giá trị $w$ tương ứng với một bài toán RL đơn mục tiêu chuẩn, agent
PPO sẽ hội tụ về một chính sách $\pi_w^*$ tối ưu cục bộ tại điểm vận
hành đó. Bằng cách \textbf{quét trọng số}
$w \in \{0{,}1;\, 0{,}2;\, \dots;\, 0{,}9\}$ và huấn luyện độc lập chín
agent, đồ án thu được chín điểm xấp xỉ trên Pareto front. Phương pháp
này có hai lợi thế: (i)~tận dụng được mọi thuật toán RL đơn mục tiêu
sẵn có; (ii)~dễ song song hoá.

\textbf{Hạn chế.} Linear scalarization không thể khôi phục được những
đoạn \textit{lõm} (concave) của Pareto front. Đây là điều cần lưu ý khi
diễn giải kết quả ở Phase 2.

\subsection{Action Masking}
Trong môi trường ELDAS, không phải hành động nào cũng hợp lệ tại mỗi
trạng thái: nếu host $h_j$ không đủ tài nguyên cho tác vụ $\tau_i$, hành
động ``đặt $\tau_i$ lên $h_j$'' phải bị cấm. Cách xử lý phổ biến trong
RL là \textbf{action masking}~\cite{huang2022maskable}: ở mỗi bước, môi
trường trả về kèm theo một danh sách các host khả thi, và những host
không hợp lệ bị loại khỏi phân phối xác suất \emph{trước khi} agent
chọn hành động --- nói cách khác, xác suất chọn một host không đủ tài
nguyên bằng đúng $0$. ELDAS sử dụng \textit{MaskablePPO} trong
\texttt{sb3-contrib} --- thư viện đã hiện thực sẵn cơ chế này. Cách
làm đem lại hai lợi ích: (i)~loại bỏ hoàn toàn khả năng vi phạm ràng
buộc capacity; (ii)~agent học nhanh hơn nhờ không phải lãng phí mẫu
vào các hành động sẽ bị phạt cứng.

% ============================================================
\section{Nghiên cứu liên quan}
\label{sec:related-work}

Phần này khảo sát năm công trình tiêu biểu, thuộc ba hướng nghiên cứu
liên quan trực tiếp đến ELDAS: RL cho lập lịch tài nguyên, lập lịch đa
mục tiêu kèm yếu tố năng lượng, và các framework mô phỏng / MORL.

\subsection{DeepRM (Mao và cộng sự, 2016)}
DeepRM~\cite{mao2016resource} là một trong những công trình đầu tiên áp
dụng \textit{Deep Reinforcement Learning} cho bài toán cấp phát tài
nguyên cluster. Đóng góp chính của bài báo gồm: (i)~biểu diễn trạng
thái cluster bằng \textit{ảnh hai chiều} (thời gian $\times$ tài
nguyên), từ đó tận dụng được kiến trúc CNN; (ii)~chỉ ra rằng agent DRL
có thể hội tụ và vượt heuristic Shortest-Job-First trong môi trường mô
phỏng. \textbf{Hạn chế:} chỉ tối thiểu một mục tiêu duy nhất là
\textit{average job slowdown}, không xét đến năng lượng và chỉ thử
nghiệm ở quy mô nhỏ ($\le 20$ tác vụ đồng thời).

\subsection{Decima (Mao và cộng sự, 2019)}
Decima~\cite{decima2019} mở rộng hướng đi của DeepRM, dùng \textit{Graph
Neural Network} (GNN) để mô hình hoá Directed Acyclic Graph của các
Spark job. Hệ thống đã được triển khai thực tế trên cụm $25$ máy và
giảm \textit{average job completion time} $21\%$ so với bộ lập lịch
mặc định của Spark. \textbf{Liên hệ với ELDAS:} cùng triết lý ``RL +
simulation + deployment''; việc đưa GNN vào ELDAS được xếp ở mục
\textit{stretch goal} trong Phase 2, sẽ chỉ tiến hành sau khi MLP đã
hội tụ.

\subsection{NSGA-II cho điều phối đám mây tiết kiệm năng lượng}
Một hướng nghiên cứu khác dùng \textit{Non-dominated Sorting Genetic
Algorithm II} (NSGA-II)~\cite{deb2002fast} để giải bài toán đa mục tiêu
năng lượng + QoS. Ưu điểm của hướng này là thuật toán tiến hoá có thể
sinh ra cả Pareto front trong một lần chạy mà không cần quét trọng
số. Tuy nhiên, mỗi lần ra quyết định lại đòi hỏi đánh giá hàng nghìn
cá thể, nên không phù hợp với \textit{online scheduling} thời gian
thực. Đây cũng chính là khoảng trống mà RL có lợi thế: chính sách đã
huấn luyện chỉ cần một bước forward pass duy nhất khi suy luận.

\subsection{Beware of Fragmentation (Weng và cộng sự, 2023)}
Công trình~\cite{weng2023alibaba} công bố Alibaba PAI Cluster Trace
v2023 và đề xuất \textit{Fragmentation Gradient Descent} --- một
heuristic giúp giảm phân mảnh GPU tới $10\%$. \textbf{Liên hệ trực
tiếp với ELDAS:} đây chính là nguồn dữ liệu workload mà ELDAS sử dụng,
đồng thời bài báo cũng cung cấp một baseline heuristic đủ mạnh để so
sánh ở Phase 2.

\subsection{MO-Gymnasium (Alegre và cộng sự, 2022)}
MO-Gymnasium~\cite{alegre2022mogym} là framework chuẩn cho các
benchmark MORL, mở rộng API \texttt{step()} của Gymnasium để trả về
vector phần thưởng. Thư viện cung cấp một số môi trường tham chiếu
(\textit{deep-sea-treasure}, \textit{minecart}, \textit{four-room})
cùng các chỉ số đánh giá Pareto front (hypervolume, sparsity). ELDAS
đăng ký một môi trường custom theo đúng chuẩn này, từ đó có thể
tương thích với mọi thuật toán MORL trong các nghiên cứu sau.

\begin{table}[H]
  \centering
  \caption{Bảng tổng hợp so sánh các nghiên cứu liên quan.}
  \label{tab:related-work}
  \footnotesize
  \begin{tabular}{lcccccc}
    \toprule
    \textbf{Nghiên cứu} & \textbf{Năm} & \textbf{Thuật toán}
                       & \textbf{Đa mục tiêu?} & \textbf{Năng lượng?}
                       & \textbf{Simulator} & \textbf{GPU?} \\
    \midrule
    DeepRM      & 2016 & DQN+CNN  & Không & Không & Custom    & Không \\
    Decima      & 2019 & GNN+RL   & Không & Không & Spark     & Không \\
    NSGA-II/EE  & 2020 & GA       & Có    & Có    & CloudSim  & Không \\
    FGD/Alibaba & 2023 & Heuristic& Có    & Không & Trace     & Có    \\
    MO-Gym      & 2022 & Framework& Có    & ---   & Gym       & ---   \\
    \textbf{ELDAS} & \textbf{2026} & \textbf{MORL+PPO}
                & \textbf{Có} & \textbf{Có} & \textbf{CloudSim Plus}
                & \textbf{Có} \\
    \bottomrule
  \end{tabular}
\end{table}

\textbf{Định vị đóng góp.} Bảng~\ref{tab:related-work} cho thấy ELDAS
lấp vào khoảng giao của bốn yêu cầu: \textit{đa mục tiêu}, \textit{tích
hợp năng lượng}, \textit{theo dõi GPU} và \textit{sử dụng workload
trace thực}, tất cả trong cùng một nền tảng MORL. Theo tìm hiểu của
tác giả tại thời điểm viết báo cáo (2026), chưa có công trình mã nguồn
mở nào kết hợp đồng thời cả bốn đặc tính này.

% ============================================================
\section{Công nghệ nền tảng}
\label{sec:cong-nghe-nen-tang}

\subsection{CloudSim Plus}
CloudSim Plus~\cite{filho2017cloudsim} là framework mô phỏng \textit{rời
rạc sự kiện} (\textit{discrete-event simulation}) cho điện toán đám
mây, viết bằng Java 21. So với CloudSim gốc (Đại học Melbourne, 2009),
CloudSim Plus được tái cấu trúc lại theo các nguyên tắc kỹ thuật phần
mềm hiện đại: API dạng \textit{fluent}, tách bạch rõ ràng
\texttt{Host} / \texttt{Vm} / \texttt{Cloudlet}, và tích hợp sẵn
\texttt{PowerModel}.

\textbf{Lý do chọn:} (i)~mã nguồn mở, cộng đồng còn duy trì tích cực;
(ii)~hoạt động \textit{event-driven}, cho phép tạm dừng và tiếp tục mô
phỏng --- đây là yêu cầu bắt buộc của vòng lặp RL (xem §4.5);
(iii)~đã có sẵn mô hình năng lượng cho CPU. \textbf{Hạn chế:} chưa hỗ
trợ GPU, do đó ELDAS bổ sung thêm lớp tracking riêng.

\paragraph{Bản chất mô phỏng host: POJO + DES, không Thread/Process.}
Khác với một số công cụ mô phỏng đời đầu (SimJava và các nguyên mẫu
CloudSim thuở sơ khai) vốn ánh xạ mỗi thực thể mô phỏng
(Datacenter, Broker, Host, VM) thành một \texttt{java.lang.Thread}
riêng biệt, CloudSim Plus
\textbf{tái thiết kế toàn bộ theo mô hình hướng đối tượng điều khiển
bằng sự kiện} (\textit{event-driven object model}). Mọi
\texttt{Host}, \texttt{Vm}, \texttt{Cloudlet} thực chất chỉ là các
\textit{Plain Old Java Object} (POJO) cư trú trên Heap của JVM,
\emph{không} phải là Thread hay tiến trình hệ điều hành. Toàn bộ vòng
lặp mô phỏng chạy trên \textbf{một luồng thực thi duy nhất}, liên tục
rút đối tượng \texttt{SimEvent} từ một \textit{Priority Queue} có tên
\texttt{FutureQueue} (sắp xếp theo timestamp), chuyển sang
\texttt{DeferredQueue} khi đã đến hạn, rồi gọi trực tiếp phương thức
\texttt{processEvent()} của đối tượng đích. Đồng hồ mô phỏng
(\textit{simulation clock}) không trôi theo thời gian thực mà
\emph{nhảy vọt} từ sự kiện này sang sự kiện kế tiếp, nhờ vậy có thể
``tua nhanh'' qua các khoảng trống không có sự kiện. Quyết định kiến
trúc này quan trọng vì ba lý do: (i)~triệt tiêu chi phí
\textit{context switch} và \textit{lock contention} của hệ điều hành
khi quy mô lên tới hàng vạn host; (ii)~đảm bảo \textbf{tính xác định}
(determinism) --- cùng một random seed luôn cho ra kết quả giống hệt
nhau, điều kiện tiên quyết cho nghiên cứu khoa học; (iii)~tiết kiệm
bộ nhớ vì mỗi Thread JVM tốn $\sim 512$\,KB stack, trong khi một POJO
chỉ tốn vài chục byte.

\paragraph{Ưu điểm.} (i)~\textit{Reproducibility tuyệt đối} nhờ
single-threaded loop và POJO; (ii)~API \textit{fluent} và mô hình
object-chaining (\texttt{cloudlet.getVm().getHost()}) giảm boilerplate;
(iii)~hỗ trợ sẵn \texttt{PowerModel}, các bộ lập lịch CPU
(\texttt{TimeShared}, \texttt{SpaceShared}, mô phỏng CFS của Linux),
auto-scaling theo \textit{Event Listener}, và đọc trace từ Google
Cluster / Alibaba dưới dạng stream; (iv)~có cơ chế \texttt{startSync()}
cho phép \emph{tạm dừng} mô phỏng tại mỗi khoảng \texttt{runFor()} ---
đúng cơ chế ELDAS dùng để chen vòng lặp RL qua Py4J (§4.5);
(v)~thực nghiệm của nhóm phát triển trên kịch bản $20\,000$ host,
$40\,000$ VM, $50\,000$ Cloudlet cho thấy CloudSim Plus chạy nhanh hơn
CloudSim 4.0.0 từ $20\%$ tới hơn $90\%$ tuỳ thuật toán cấp
phát~\cite{cloudsimplus2024perf}.

\paragraph{Hạn chế và điểm nghẽn ở quy mô lớn.} Ngược lại, kiến trúc
này cũng kéo theo bốn ràng buộc cần lưu ý khi xác định phạm vi đồ án:
\begin{enumerate}
  \item \textbf{Lõi đơn luồng không tận dụng được nhiều core.} Cho
        phép chạy \emph{song song nhiều kịch bản độc lập} (mỗi kịch
        bản một JVM hoặc một luồng lõi riêng), nhưng \emph{một}
        Datacenter không thể phân tán event-loop ra nhiều CPU. Trần
        thông lượng bị chặn bởi single-thread performance của CPU.
  \item \textbf{Độ phức tạp $O(V \cdot H)$ của allocation policy.}
        Các chính sách scoring đầy đủ (BestFit, WorstFit, K8s LRP,
        meta-heuristic) phải duyệt qua mọi host khả thi cho mỗi VM
        cần đặt; khi $V, H$ cùng đạt $10^4$, khâu cấp phát có thể tốn
        hàng giờ CPU.
  \item \textbf{Áp lực \textit{recursive update}.} Mỗi nhịp
        \texttt{schedulingInterval}, Datacenter gọi đệ quy\\
        \texttt{updateVmsProcessing()} qua toàn bộ Host $\to$ VM $\to$
        Cloudlet. Thí nghiệm trên Google Cluster Trace cho thấy với
        $10\,000$ host, mô phỏng \emph{hai ngày} dữ liệu thực có thể
        tốn tới \emph{bảy ngày} wall-time;
        điểm nghẽn lớn nhất là vòng lặp cập nhật chứ không phải
        thuật toán cấp phát.
  \item \textbf{Áp lực JVM Garbage Collector.} Một episode dài với
        hàng triệu \texttt{SimEvent} và\\ \texttt{HostStateHistoryEntry}
        làm Heap đầy nhanh; các đợt \textit{Stop-The-World GC} của
        G1GC/Shenandoah có thể đóng băng tiến trình hàng giây. Trần
        thực tế của một JVM $16$\,GB RAM nằm vào khoảng
        $10^4$--$10^5$ host trên cùng một mô phỏng.
  \item \textbf{Không có GUI.} Cấu hình hoàn toàn bằng mã Java hoặc
        YAML; rào cản đối với người dùng không lập trình hệ thống.
\end{enumerate}

\paragraph{Định vị quy mô datacenter của ELDAS.}
\label{para:datacenter-scale}
Việc thấu hiểu các điểm nghẽn nói trên trực tiếp ấn định
\emph{khoảng quy mô khả thi} của một thí nghiệm CloudSim Plus, và
qua đó định vị ELDAS trên thang đo công nghiệp.
Bảng~\ref{tab:datacenter-scale} so sánh ba mức quy mô tham chiếu.

\begin{table}[H]
  \centering
  \caption{So sánh quy mô datacenter ELDAS với các mức công nghiệp.}
  \label{tab:datacenter-scale}
  \footnotesize
  \begin{tabular}{lccc}
    \toprule
    \textbf{Tiêu chí} & \textbf{ELDAS (Phase 1)} &
        \textbf{Normal-scale} & \textbf{Hyperscale} \\
    \midrule
    Số host vật lý        & $10$ (mở rộng được $50$)
                          & $10^2$--$10^3$
                          & $10^5$--$10^6$ \\
    Loại đại diện         & Lab / micro-cluster
                          & Cluster học thuật, edge DC
                          & Google, AWS, Azure region \\
    \#vCPU/host           & $16$ (cấu hình được)
                          & $32$--$128$
                          & $96$--$192$ \\
    GPU/host              & $0$--$4$ (mô phỏng)
                          & $0$--$8$
                          & $8$+ (cụm AI dedicated) \\
    \#task/episode        & $\sim 2 \cdot 10^3$ (Alibaba)
                          & $10^4$--$10^5$
                          & $10^7$+ trên ngày \\
    PUE điển hình         & $1{,}5$ (giả định)
                          & $1{,}5$--$1{,}8$
                          & $1{,}10$--$1{,}15$~\cite{google2023datacenter} \\
    Khả năng tái lập      & Tuyệt đối (1 seed)
                          & Khó (multi-tenant)
                          & Không thể \\
    \bottomrule
  \end{tabular}
\end{table}

ELDAS rõ ràng thuộc nhóm \textbf{lab-scale / micro-cluster}: $10$
host, vài nghìn task/episode trên trace Alibaba đã được lọc về một
window đại diện. Mức quy mô này có hai hệ quả cần khẳng định trong
phần giới hạn:
\begin{itemize}
  \item \textbf{Đủ để chứng minh thuật toán.} Với $10$ host và
        $6$ scheduler, không gian hành động (10 host) và không
        gian trạng thái ($6H+4 = 64$ chiều) đủ rộng để tạo ra
        \emph{Pareto front có thể quan sát}: chênh lệch năng lượng
        BestFit$\leftrightarrow$RoundRobin đạt $\sim 20$--$26\%$ (xem
        Chương~5), và PPO-min đã có vị trí interior Pareto thực sự
        --- đủ chứng minh giả thuyết MORL khả thi trên CloudSim
        Plus.
  \item \textbf{Không đại diện cho hyperscale.} Tại quy mô
        $10^5$+ host, các yếu tố trở nên chi phối mà ELDAS không
        mô phỏng: phân tầng mạng đa lớp, fault domain, multi-zone
        replication, năng lượng làm mát phụ thuộc nhiệt độ ngoài
        trời, và đặc biệt là chi phí lập lịch đệ quy phình theo
        $O(V \cdot H)$ đã nêu. Một thí nghiệm CloudSim Plus
        $10\,000$+ host sẽ vướng cả ba điểm nghẽn ở mục trên và đòi
        hỏi hạ tầng JVM riêng ($\geq 32$\,GB Heap, GC tinh chỉnh).
        Đây là giới hạn được nêu rõ trong Chương~\ref{chap:ket-luan}.
\end{itemize}

\subsection{Py4J -- cầu nối Java--Python}
Py4J~\cite{py4j2024} là thư viện cho phép code Python gọi trực tiếp các đối
tượng Java trong cùng JVM thông qua \textit{GatewayServer}. Cơ chế:
\begin{itemize}
  \item Phía Java khởi tạo \texttt{GatewayServer} lắng nghe trên TCP port
        $25333$;
  \item Phía Python tạo \texttt{JavaGateway()} để truy cập đối tượng
        \texttt{entry\_point};
  \item Mỗi lời gọi method được serialize qua TCP socket, kết quả trả về
        đồng bộ.
\end{itemize}

\begin{table}[H]
  \centering
  \caption{So sánh các tùy chọn cầu nối Java--Python cho ELDAS.}
  \label{tab:bridge-comparison}
  \small
  \begin{tabular}{lccc}
    \toprule
    \textbf{Tiêu chí} & \textbf{Py4J} & \textbf{gRPC} & \textbf{REST} \\
    \midrule
    Độ trễ một call    & $\sim 0{,}1$~ms & $\sim 1$~ms     & $\sim 5$~ms     \\
    Cần định nghĩa schema? & Không   & Có (\texttt{.proto}) & Có (OpenAPI) \\
    Truyền đối tượng phức tạp & Trực tiếp & Phải serialize & JSON \\
    Phù hợp tần suất cao? & Tốt   & Khá          & Kém       \\
    \bottomrule
  \end{tabular}
\end{table}

Do vòng lặp RL có thể gọi \texttt{step()} tới hàng nghìn lần mỗi
episode, Py4J là phương án cân bằng tốt nhất giữa hiệu năng và độ phức
tạp triển khai.

\subsection{MO-Gymnasium}
MO-Gymnasium là một \textit{fork} của Gymnasium do Alegre và cộng sự
duy trì~\cite{alegre2022mogym}. Khác biệt quan trọng nhất nằm ở API
\texttt{step(action)}: hàm này trả về \texttt{(obs, reward\_vector,
terminated, truncated, info)}, với \texttt{reward\_vector} $\in
\mathbb{R}^d$ thay vì một số thực. ELDAS đăng ký môi trường
\texttt{ELDAS-v0} kế thừa từ lớp \texttt{MOEnv}, qua đó có thể
\textit{plug-and-play} với mọi thuật toán MORL tuân theo chuẩn này.

\subsection{Docker Compose}
Toàn bộ hệ thống được đóng gói bằng Docker Compose v2, gồm hai service
(\texttt{cloudsim-java}, \texttt{rl-agent}) cùng ba volume chia sẻ
(\texttt{trace-data}, \texttt{metrics}, \texttt{checkpoints}). Mục tiêu
là khả năng \textbf{tái lập kết quả} (\textit{reproducibility}): bất
kỳ ai có Docker đều có thể chạy lại thực nghiệm bằng một câu lệnh duy
nhất:
\begin{lstlisting}[style=yamlstyle, caption={Khởi chạy hệ thống.}, label={lst:run}]
docker compose up --build
\end{lstlisting}

\subsection{So sánh với KWOK}
KWOK (\textit{Kubernetes WithOut Kubelet}) là một lựa chọn thay thế khá
phổ biến để mô phỏng cluster K8s ở quy mô lớn.
Bảng~\ref{tab:kwok-vs-cloudsim} tóm lược lý do ELDAS chọn CloudSim Plus
thay vì KWOK ở Phase 1.

\begin{table}[H]
  \centering
  \caption{So sánh KWOK và CloudSim Plus theo các tiêu chí của ELDAS.}
  \label{tab:kwok-vs-cloudsim}
  \small
  \begin{tabular}{lcc}
    \toprule
    \textbf{Tiêu chí} & \textbf{KWOK} & \textbf{CloudSim Plus} \\
    \midrule
    Mô hình thời gian       & Real-time         & Discrete-event \\
    Mô hình năng lượng      & Không có          & Có sẵn (PowerModel) \\
    Pause/Resume mô phỏng   & Khó               & Hỗ trợ nguyên thủy \\
    Quy mô tối đa           & Hàng nghìn node   & Hàng trăm host \\
    Phù hợp huấn luyện RL?  & Khó (không pause) & Tốt \\
    Hỗ trợ GPU              & Plugin            & Phải bổ sung \\
    \bottomrule
  \end{tabular}
\end{table}

Lợi thế quyết định của CloudSim Plus nằm ở cơ chế
\textit{discrete-event}: mô phỏng có thể ``tua nhanh'' qua các khoảng
không có sự kiện, nhờ vậy một episode RL chỉ mất vài giây thay vì vài
giờ thời gian thực.

% ============================================================
\section{Kết luận chương}
\label{sec:ket-luan-c2}

Chương này đã đặt nền tảng lý thuyết cho ba quyết định thiết kế cốt lõi
của ELDAS: (i)~chọn \textbf{mô hình năng lượng tuyến tính} vì cân bằng
được độ chính xác và khả năng tái lập; (ii)~chọn \textbf{PPO kết hợp
linear scalarization} làm thuật toán nền cho Phase 2 nhờ tính ổn định
và dễ song song hoá; (iii)~chọn ngăn xếp công nghệ \textbf{CloudSim
Plus + Py4J + MO-Gymnasium} vì khả năng pause/resume mô phỏng và mức
độ tương thích với chuẩn MORL. Bảng tổng hợp ở §2.4 cũng đã định rõ vị
trí và đóng góp của ELDAS trong tương quan với các công trình liên
quan. Chương kế tiếp sẽ trình bày kiến trúc bốn lớp hiện thực hoá các
quyết định nói trên.
